% Generated draft: requires mathematical and source review.
\documentclass[11pt]{article}
\usepackage[margin=25mm]{geometry}
\usepackage{fontspec}
\setmainfont{Latin Modern Roman}
\usepackage{amsmath,amssymb,mathtools,unicode-math}
\setmathfont{Latin Modern Math}
\usepackage{enumitem,needspace,xurl,hyperref}
\setlength{\parindent}{0pt}
\setlength{\parskip}{5pt}
\setlength{\emergencystretch}{3em}
\newcommand{\N}{\mathbb{N}}
\newcommand{\Z}{\mathbb{Z}}
\newcommand{\Q}{\mathbb{Q}}
\newcommand{\R}{\mathbb{R}}
\newcommand{\C}{\mathbb{C}}
\newcommand{\F}{\mathbb{F}}
\newcommand{\A}{\mathbb{A}}
\newcommand{\PP}{\mathbb{P}}
\newcommand{\E}{\mathbb{E}}
\newcommand{\eps}{\varepsilon}
\newcommand{\dd}{\,\mathrm{d}}
\begin{document}
\section*{UnsolvedMath ID: 20002016}
\textbf{Source problem code:} catalogue entry
\subsection{Definitions and mathematical statement}
This extraction combines background on form analogues of \(Q\)-curvature with a nonspecific ``Problems \(k\)'' program. Work on a closed oriented Riemannian manifold \(M\) of even dimension \(n\), not necessarily connected. Write \(\Omega^k\) for smooth real forms, \(C^k=\ker d\), and use de Rham cohomology. For a conformal change \(\widehat g=e^{2\omega}g\), the scalar \(Q\)-curvature is best recorded as its top-form density \(q_g\). The source gives
\[
 q_{\widehat g}=q_g+P_n\omega,
\]
where the critical GJMS operator \(P_n:\Omega^0\to\Omega^n\) is formally self-adjoint and has a detour factorization of type \(dM_1d\), up to conventions. Consequently multiplication by \(q_g\) on locally constant functions induces a conformally invariant map \(C^0\to H^n(M;\R)\), and \(\int_M u\,q_gc\) is invariant for \(c\in C^0\), \(u\in\ker P_n\). Its value on \(1\) is \(q_g\). One may alter an operator extension off locally constant functions by terms annihilating them, while the source asks that formal self-adjointness be retained.

The general reported construction has \(M_k^g:\Omega^k\to\Omega^{n-k}\), \(0\leq k\leq n/2-1\), obtained using ambient metrics and tractors. In the form/density identifications intended by the source, its axioms are
\[
 M_k^{\widehat g}c=M_k^gc+L_k(\omega c)\quad(c\in C^k),\qquad
 L_k=\beta_k dM_{k+1}^gd,
\]
with fixed normalization constants, and formal self-adjointness for both \(M_k^g\) and \(L_k\). The terminal factor at the middle degree requires the construction's endpoint convention, not supplied here. Put
\[
 \mathcal H_g^k=\{c\in C^k:dM_k^gc=0\},\qquad
 \Phi_g(c)=[M_k^gc]\in H^{n-k}(M;\R).
\]
The source states conformal invariance of \(\mathcal H_g^k\) and \(\Phi_g\), nontriviality on some conformal manifolds, and the invariant pairing
\[
 \int_M u\wedge M_k^gc\quad(c\in C^k,\ u\in\ker L_k),
\]
with signs/weighted identifications absorbed in the formal-adjoint convention. Also \(M_0^g1=q_g\). These are reported established form-\(Q\) properties, not hypotheses already verified for every arbitrary operator.

The aid's minimal-axiom conclusion is that the transformation and detour laws alone make \(dM_k^g c\) independent of conformal representative on closed inputs, and the change in \(M_k^g c\) exact. Formal self-adjointness gives the kernel-pairing invariance. A separate cohomological descent issue remains. Let
\[
 q:\mathcal H^k\to H^k(M;\R),\qquad c\mapsto[c].
\]
Then \(\Phi\) factors through the represented subspace \(q(\mathcal H^k)\) if and only if
\[
 [M_k^g(d\alpha)]=0
 \quad\text{for every }d\alpha\in\mathcal H^k.
\]
This exact-input condition is necessary and sufficient for independence of the chosen conformal-harmonic representative. Obtaining the resulting map on all of \(H^k(M;\R)\) additionally requires surjectivity of \(q\). Neither conformal invariance nor equality of a dimension with a Betti number automatically supplies both gates.

The source's near-middle illustration, for \(n\geq4\) and \(p=n/2-1\), uses the weighted operator
\[
 \mathcal M_g=d\delta+2J-4P^{\sharp}:\Omega^p\to\Omega^p[-2],
 \qquad J=\frac{\operatorname{Scal}}{2(n-1)},\quad
 P=\frac{\operatorname{Ric}-Jg}{n-2},
\]
where \(P^{\sharp}\) acts on form slots. The weighted target identifies by conformal Hodge star with \((p+2)\)-forms. Exercise \(7\) asks to check on closed inputs the law
\[
 \mathcal M_{\widehat g}c=\mathcal M_gc+\beta\delta d(\omega c),\qquad\beta\ne0,
\]
under the intended weighted identification. It notes formal self-adjointness of the Maxwell operator \(\delta d\), invariant pairings with its kernel, invariance modulo \(\operatorname{im}\delta\), and invariance of \(\delta\mathcal M_g\) on closed forms. On a closed Riemannian manifold \(\ker(\delta d)\) consists of closed forms. The source also reports a conformally invariant \(\delta\mathcal M_gd\) on \((n/2-2)\)-forms, giving Paneitz in dimension four, and defines the corresponding conformal harmonics by \(\delta\mathcal M_gc=0\).

It cautions that their dimension need not always equal the Betti number, suggests a generic equality through elliptic coercivity of \((d,\delta\mathcal M)\) without defining a small exceptional set, and gives \(S^p\times S^{p+2}\) with standard product metric as a nontrivial example where conformal harmonics in degree \(p\) are ordinary harmonic forms. It further asserts a non-removability-by-invariant-operator-and-divergence property, referring to Maxwell leading-order uniqueness and an unavailable formula \((44)\). These background claims need the precise conventions and references.

The final ``Problems \(k\)'' only asks for analogues of unspecified \(Q\)-curvature conundrums for the \(M_k^g\). The initial ``Problem \(2\) have a role'' fragment and the following Chapter C heading with Thomas Branson's ``Anti-conformal perturbations'' are boundary material, not complete additional problems. No characterization of \(Q\), general cohomology-surjectivity theorem, or precise new open-program statement is invented from them.
\subsection{Short English statement}
The form-\(Q\) construction gives conformally invariant harmonic-input spaces, output cohomology classes and kernel pairings. The aid isolates two separate requirements for a map on ordinary de Rham cohomology: exact inputs must map to exact outputs, and every cohomology class must have an admissible representative. The source's remaining open prompt is nonspecific.
\subsection{Sources}
\begin{itemize}
\item Local UnsolvedMath registry, permanent ID \(20002016\), source code \texttt{catalogue entry}; canonical collection \texttt{aim-geometry-notes.json}, notes[353]. The source-item field is \(2\), but the extraction covers the end of Chapter B, equations \((35),(36)\), Exercise \(7\), and ``Problems \(k\).''
\item AIM workshop ``Conformal Structure in Geometry, Analysis, and Physics,'' 12--16 August 2003, Palo Alto, California: \url{https://aimath.org/WWN/confstruct/confstruct.pdf}. The supplied verification reports comparison of official HTML and PDF. Page numbers \(22,23\), the Chapter C heading and Thomas Branson's next-topic attribution are boundary material. No external source or PDF was opened in this session.
\item The aid attributes the central established construction to Branson and Gover; the source names \(Q\)-curvature, GJMS operators, Maxwell, Paneitz, ambient metrics and tractors. No full authoritative bibliographic entry, formula \((44)\), or endpoint/normalization conventions are supplied. Original remarks and literature fields are empty, with no proposer or year.
\item Supplied research attempt \(1\), \texttt{partial\_result}, run status \texttt{partially\_solved}, rigor label ``proved,'' novelty confidence ``low.'' Registry classification \texttt{PARTIAL-PROGRESS}, updated 14 August 2026. Its candidate contribution is the two-gate descent criterion, not a claim to invent the form-\(Q\) construction.
\end{itemize}
\end{document}
